% GENERATED FILE. Edit canonical lessons, then run npm run generate:books.
% canonical-source-hash: fnv1a64:a5c0111ecac5e815
% canonical-lessons: LA-C404-rhetor, LA-C404-rhetorica, LA-C404-dialectica, LA-C404-geometria, LA-C404-elementum, LA-R404-review-397-404

\chapter{What Pupils Studied}
\label{ch:la-what-pupils-studied}
\hlchaptermodality{404}

\begin{chapteropening}
\textbf{By the end of this chapter:} \emph{I can name the teacher of speaking, rhetoric, logic and geometry, and the first elements of a subject.}

Complete the last lesson of chapter 404: I can name the teacher of speaking, rhetoric, logic and geometry, and the first elements of a subject.
\end{chapteropening}

\section[rhētor, rhētoris]{rhētor, rhētoris — a teacher of public speaking}

\label{lesson:LA-C404-rhetor}



\begin{quote}
Before the new one: say the Latin for knowing, then the Latin for inexperienced.
\end{quote}

\subsection*{You'll want to know: rhētor, rhētoris}
\textbf{rhētor, rhētoris} — \textquotedblleft{}a teacher of public speaking\textquotedblright{}.

\begin{quote}
Iuvenēs apud rhētorem dīcere discunt.
\end{quote}

\begin{quote}
\emph{Young men learn to speak in public with a rhetoric teacher.}
\end{quote}

\emph{rhētor, rhētoris} m., borrowed from Greek \emph{rhētōr}, a speaker. The last stage of a Roman education was the school of the \emph{rhētor}.

\begin{cousinweb}
\emph{rhētor} was borrowed from Greek \emph{rhētōr} \textquotedblleft{}a public speaker\textquotedblright{} (EtymDB; Lewis \& Short). English \emph{rhetor} and \emph{rhetoric} come from the same Greek word.
\end{cousinweb}

\begin{grammarlens}[title={What you've built}]
One more word for telling, explaining and describing at B1.
\end{grammarlens}

\subsection*{Your turn}
\begin{itemize}
\raggedright
  \item \emph{Say it:} \emph{rhētor, rhētoris}
  \item \emph{Say it:} \emph{rhētor, rhētoris}, once more
  \item \emph{Say it:} a sentence of your own with \emph{rhētor, rhētoris}
  \item \emph{Recall:} say the Latin for knowing, then the Latin for inexperienced, then say \emph{rhētor, rhētoris} again
\end{itemize}

\begin{tcolorbox}[breakable,colback=teal!4,colframe=teal!35!black,title={Before you move on}]
What does rhētor, rhētoris mean? (\textquotedblleft{}A teacher of public speaking\textquotedblright{}.) Say it once more.
\end{tcolorbox}
\section[rhētorica, rhētoricae]{rhētorica, rhētoricae — rhetoric, the art of speaking}

\label{lesson:LA-C404-rhetorica}



\begin{quote}
Before the new one: say the Latin for inexperienced, then the Latin for a teacher of public speaking.
\end{quote}

\subsection*{You'll want to know: rhētorica, rhētoricae}
\textbf{rhētorica, rhētoricae} — \textquotedblleft{}rhetoric, the art of speaking\textquotedblright{}.

\begin{quote}
Rhētoricam in scholā didicit.
\end{quote}

\begin{quote}
\emph{He learned rhetoric at school.}
\end{quote}

\emph{rhētorica, rhētoricae} f., also \emph{ars rhētorica}. English \emph{rhetoric}.

\begin{cousinweb}
\emph{rhētorica} was borrowed from Greek \emph{rhētorikē} \textquotedblleft{}the art of speaking\textquotedblright{} (EtymDB), built on \emph{rhētōr} \textquotedblleft{}a speaker\textquotedblright{}. English \emph{rhetoric} comes from it.
\end{cousinweb}

\begin{grammarlens}[title={What you've built}]
One more word for telling, explaining and describing at B1.
\end{grammarlens}

\subsection*{Your turn}
\begin{itemize}
\raggedright
  \item \emph{Say it:} \emph{rhētorica, rhētoricae}
  \item \emph{Say it:} \emph{rhētorica, rhētoricae}, once more
  \item \emph{Say it:} a sentence of your own with \emph{rhētorica, rhētoricae}
  \item \emph{Recall:} say the Latin for inexperienced, then the Latin for a teacher of public speaking, then say \emph{rhētorica, rhētoricae} again
\end{itemize}

\begin{tcolorbox}[breakable,colback=teal!4,colframe=teal!35!black,title={Before you move on}]
What does rhētorica, rhētoricae mean? (\textquotedblleft{}Rhetoric, the art of speaking\textquotedblright{}.) Say it once more.
\end{tcolorbox}
\section[dialectica, dialecticae]{dialectica, dialecticae — logic, dialectic}

\label{lesson:LA-C404-dialectica}



\begin{quote}
Before the new one: say the Latin for a teacher of public speaking, then the Latin for rhetoric.
\end{quote}

\subsection*{You'll want to know: dialectica, dialecticae}
\textbf{dialectica, dialecticae} — \textquotedblleft{}logic, dialectic\textquotedblright{}.

\begin{quote}
Dialectica docet vēra ā falsīs distinguere.
\end{quote}

\begin{quote}
\emph{Logic teaches us to tell true from false.}
\end{quote}

\emph{dialectica, dialecticae} f., from Greek \emph{dialektikē}, the art of discussion: the logic of the philosophers. English \emph{dialectic}.

\begin{cousinweb}
\emph{dialectica} was borrowed from Greek \emph{dialektikē} \textquotedblleft{}the art of discussion\textquotedblright{} (Lewis \& Short), built on a verb meaning \textquotedblleft{}to talk things through\textquotedblright{}. English \emph{dialectic} comes from it.
\end{cousinweb}

\begin{grammarlens}[title={What you've built}]
One more word for telling, explaining and describing at B1.
\end{grammarlens}

\subsection*{Your turn}
\begin{itemize}
\raggedright
  \item \emph{Say it:} \emph{dialectica, dialecticae}
  \item \emph{Say it:} \emph{dialectica, dialecticae}, once more
  \item \emph{Say it:} a sentence of your own with \emph{dialectica, dialecticae}
  \item \emph{Recall:} say the Latin for a teacher of public speaking, then the Latin for rhetoric, then say \emph{dialectica, dialecticae} again
\end{itemize}

\begin{tcolorbox}[breakable,colback=teal!4,colframe=teal!35!black,title={Before you move on}]
What does dialectica, dialecticae mean? (\textquotedblleft{}Logic, dialectic\textquotedblright{}.) Say it once more.
\end{tcolorbox}
\section[geōmetria, geōmetriae]{geōmetria, geōmetriae — geometry}

\label{lesson:LA-C404-geometria}



\begin{quote}
Before the new one: say the Latin for rhetoric, then the Latin for logic.
\end{quote}

\subsection*{You'll want to know: geōmetria, geōmetriae}
\textbf{geōmetria, geōmetriae} — \textquotedblleft{}geometry\textquotedblright{}.

\begin{quote}
Geōmetria terrae mēnsūram docet.
\end{quote}

\begin{quote}
\emph{Geometry teaches the measuring of land.}
\end{quote}

\emph{geōmetria, geōmetriae} f., from Greek \emph{geōmetria}, land-measuring: \emph{gē}, earth + \emph{metron}, a measure. English \emph{geometry}.

\begin{cousinweb}
\emph{geōmetria} was borrowed from Greek \emph{geōmetria} \textquotedblleft{}land-measuring\textquotedblright{} (EtymDB), built from \emph{gē} \textquotedblleft{}earth\textquotedblright{} and \emph{metron} \textquotedblleft{}a measure\textquotedblright{}. English \emph{geometry} comes from it.
\end{cousinweb}

\begin{grammarlens}[title={What you've built}]
One more word for telling, explaining and describing at B1.
\end{grammarlens}

\subsection*{Your turn}
\begin{itemize}
\raggedright
  \item \emph{Say it:} \emph{geōmetria, geōmetriae}
  \item \emph{Say it:} \emph{geōmetria, geōmetriae}, once more
  \item \emph{Say it:} a sentence of your own with \emph{geōmetria, geōmetriae}
  \item \emph{Recall:} say the Latin for rhetoric, then the Latin for logic, then say \emph{geōmetria, geōmetriae} again
\end{itemize}

\begin{tcolorbox}[breakable,colback=teal!4,colframe=teal!35!black,title={Before you move on}]
What does geōmetria, geōmetriae mean? (\textquotedblleft{}Geometry\textquotedblright{}.) Say it once more.
\end{tcolorbox}
\section[elementum, elementī]{elementum, elementī — a first principle, an element; (plural) the ABC}

\label{lesson:LA-C404-elementum}



\begin{quote}
Before the new one: say the Latin for logic, then the Latin for geometry.
\end{quote}

\subsection*{You'll want to know: elementum, elementī}
\textbf{elementum, elementī} — \textquotedblleft{}a first principle, an element; (plural) the ABC\textquotedblright{}.

\begin{quote}
Puer prīma elementa discit.
\end{quote}

\begin{quote}
\emph{The boy is learning his ABC.}
\end{quote}

\emph{elementum, elementī} n. In the plural, \emph{elementa} are the letters, the first steps of any subject, and the four elements: earth, water, air and fire. English \emph{element}.

\begin{cousinweb}
The origin of \emph{elementum} is not settled: Lewis \& Short connect it with \emph{alō} \textquotedblleft{}to nourish\textquotedblright{}, and EtymDB pairs it with Greek \emph{stoicheion} \textquotedblleft{}an element\textquotedblright{}, the word it translates. English \emph{element} comes from it.
\end{cousinweb}

\begin{grammarlens}[title={What you've built}]
One more word for telling, explaining and describing at B1.
\end{grammarlens}

\subsection*{Your turn}
\begin{itemize}
\raggedright
  \item \emph{Say it:} \emph{elementum, elementī}
  \item \emph{Say it:} \emph{elementum, elementī}, once more
  \item \emph{Say it:} a sentence of your own with \emph{elementum, elementī}
  \item \emph{Recall:} say the Latin for logic, then the Latin for geometry, then say \emph{elementum, elementī} again
\end{itemize}

\begin{tcolorbox}[breakable,colback=teal!4,colframe=teal!35!black,title={Before you move on}]
What does elementum, elementī mean? (\textquotedblleft{}A first principle, an element; (plural) the ABC\textquotedblright{}.) Say it once more.
\end{tcolorbox}
\section[(dialogue)]{Words from chapters 397–404 — a review}

\label{lesson:LA-R404-review-397-404}



\begin{quote}
Say the Latin for an oath and an agreement.
\end{quote}

\subsection*{Your turn}
Say these aloud:

\begin{itemize}
\raggedright
  \item an oath, an agreement, to strike a bargain, a pact, to bargain for
  \item to violate, to make good, to oblige, fixed, void
  \item bad faith, treacherous, a deceit, to deceive, to call to witness
  \item to teach, to educate, to train up, to equip, to instruct
  \item a slave who takes a boy to school, a teacher, a teacher of literature and grammar, an elementary teacher of letters, a schoolfellow
  \item to learn by heart, to learn thoroughly, to be eager, to think over, to call to mind
  \item learned, unlearned, educated, knowing, inexperienced
  \item a teacher of public speaking, rhetoric, logic, geometry, a first principle
\end{itemize}

\begin{tcolorbox}[breakable,colback=teal!4,colframe=teal!35!black,title={Before you move on}]
An oath? (\textbf{iūsiūrandum, iūrisiūrandī}.) An agreement? (\textbf{cōnventiō, cōnventiōnis}.) To strike a bargain? (\textbf{pacīscor, pacīscī}.) A pact? (\textbf{pactum, pactī}.) To bargain for? (\textbf{stipulor, stipulārī}.) To violate? (\textbf{violō, violāre}.) To make good? (\textbf{praestō, praestāre}.) To oblige? (\textbf{obligō, obligāre}.) Fixed? (\textbf{ratus, rata, ratum}.) Void? (\textbf{irritus, irrita, irritum}.) Bad faith? (\textbf{perfidia, perfidiae}.) Treacherous? (\textbf{perfidus, perfida, perfidum}.) A deceit? (\textbf{fallācia, fallāciae}.) To deceive? (\textbf{dēcipiō, dēcipere}.) To call to witness? (\textbf{testor, testārī}.) To teach? (\textbf{doceō, docēre}.) To educate? (\textbf{ērudiō, ērudīre}.) To train up? (\textbf{īnstituō, īnstituere}.) To equip? (\textbf{īnstruō, īnstruere}.) To instruct? (\textbf{praecipiō, praecipere}.) A slave who takes a boy to school? (\textbf{paedagōgus, paedagōgī}.) A teacher? (\textbf{praeceptor, praeceptōris}.) A teacher of literature and grammar? (\textbf{grammaticus, grammaticī}.) An elementary teacher of letters? (\textbf{litterātor, litterātōris}.) A schoolfellow? (\textbf{condiscipulus, condiscipulī}.) To learn by heart? (\textbf{ēdiscō, ēdiscere}.) To learn thoroughly? (\textbf{perdiscō, perdiscere}.) To be eager? (\textbf{studeō, studēre}.) To think over? (\textbf{meditor, meditārī}.) To call to mind? (\textbf{recordor, recordārī}.) Learned? (\textbf{doctus, docta, doctum}.) Unlearned? (\textbf{indoctus, indocta, indoctum}.) Educated? (\textbf{ērudītus, ērudīta, ērudītum}.) Knowing? (\textbf{sciēns, scientis}.) Inexperienced? (\textbf{imperītus, imperīta, imperītum}.) A teacher of public speaking? (\textbf{rhētor, rhētoris}.) Rhetoric? (\textbf{rhētorica, rhētoricae}.) Logic? (\textbf{dialectica, dialecticae}.) Geometry? (\textbf{geōmetria, geōmetriae}.) A first principle? (\textbf{elementum, elementī}.)
\end{tcolorbox}
